\par\Needspace{6\baselineskip}\medskip
\begin{center}\hypertarget{native-234-article169}{}\pdfbookmark[2]{169.}{article-169}169.\end{center}
Once the calculation in the preceding article is complete, the first task is to determine the arc $C'C''$. The most convenient way to do this is first to find, as in Article 137, the intersection $D$ of the great circles $A'C'B'$, $A''C''B''$ and their mutual inclination $\epsilon$. Then from $\epsilon$, $C'D=z'+B'D$, and $C''D=z''+B''D$, using the same formulas we gave in Article 144, we find $C'C''=v''-v'$ as well as the angles $(u',u'')$ at which the great circles $A'B'$, $A''B''$ intersect the great circle $C'C''$.

Once the arc $v''-v'$ has been found, $v'-v$ and $r$ are obtained by combining the equations
\[
r\sin(v'-v)=\frac{r''\sin(v''-v')}{P'}
\]
\[
r\sin(v'-v+v''-v')=\frac{1+P'}{P'}\cdot\frac{r'\sin(v''-v')}{1+\dfrac{Q'}{r'^3}}
\]
and similarly $r'''$ and $v'''-v''$ by combining these equations
\[
r'''\sin(v'''-v'')=\frac{r'\sin(v''-v')}{P''}
\]
\[
r'''\sin(v'''-v''+v''-v')=\frac{1+P''}{P''}\cdot\frac{r''\sin(v''-v')}{1+\dfrac{Q''}{r''^3}}\mathpunct{\cdot}
\]
\SourceLeaf{235}All the numbers found in this way would be exact if it had been possible to start with the true values of $P',P'',Q',Q''$. We would then determine the orientation of the orbital plane, as in Article 149, either from $A'C',u'$ and $\gamma'$ or from $A''C'',u''$ and $\gamma''$, and the dimensions of the orbit itself either from $r',r'',t',t''$ and $v''-v'$ or, more accurately, from $r,r''',t,t'''$ and $v'''-v$. In the first calculation, however, we shall pass over all this and concentrate on obtaining better approximations to the quantities $P',P'',Q',Q''$. We shall achieve this by applying the method explained from Article 88 onward, so that
\begin{center}
\begin{tabular}{@{}llll@{}}
from & $r,\ r',\ v'-v,\ t'-t$ & we obtain & $(\eta01)$\\
& $r',\ r'',\ v''-v',\ t''-t'$ & . . . . . . . . & $(\eta12)$\\
& $r'',\ r''',\ v'''-v'',\ t'''-t''$ & . . . . . . . . & $(\eta23)$.
\end{tabular}
\end{center}
We shall substitute these quantities and the values of $r,r',r'',r''',\cos\tfrac12(v'-v)$, etc., in formulas III—VI. This yields values of $P',Q',P'',Q''$ much more accurate than those on which the first hypothesis was based. We therefore form the second hypothesis with these values. Carrying it through in exactly the same way as the first gives still more accurate values of $P',Q',P'',Q''$ and thus leads to a third hypothesis. These operations are repeated until the values of $P',Q',P'',Q''$ appear to require no further correction. Frequent practice soon teaches one how to judge this correctly. When the heliocentric motion is small, the first hypothesis generally already gives those values with sufficient accuracy. If, however, the motion covers a larger arc, and if, in addition, the time intervals differ appreciably from one another, several successive hypotheses will be needed. In such a case the initial hypotheses require little precision in calculation. Finally, in the last hypothesis we determine the elements themselves as just indicated.
